\documentclass[11pt,a4paper]{article}

\usepackage[T1]{fontenc}
\usepackage[utf8]{inputenc}
\usepackage[portuguese]{babel}
\usepackage{lmodern}
\usepackage{microtype}
\usepackage[margin=2.5cm]{geometry}
\usepackage{xcolor}
\usepackage{graphicx}
\usepackage{booktabs}
\usepackage{tabularx}
\usepackage{enumitem}
\usepackage{tikz}
\usetikzlibrary{positioning,arrows.meta,shapes.multipart}
\usepackage[most]{tcolorbox}
\usepackage{fancyhdr}
\usepackage{titlesec}
\usepackage[hidelinks]{hyperref}

\definecolor{accent}{HTML}{1F4E79}
\definecolor{tipbg}{HTML}{EAF2FB}
\definecolor{warnbg}{HTML}{FDF1E6}
\definecolor{warnfr}{HTML}{C0611B}

\titleformat{\section}{\Large\bfseries\color{accent}}{\thesection}{1em}{}
\titleformat{\subsection}{\large\bfseries\color{accent}}{\thesubsection}{1em}{}

\setlist{itemsep=2pt,topsep=4pt}
\setlength{\parindent}{0pt}
\setlength{\parskip}{6pt}

\pagestyle{fancy}
\fancyhf{}
\fancyhead[L]{\small\color{gray}Thin Film Logbook}
\fancyhead[R]{\small\color{gray}Manual de utilização}
\fancyfoot[C]{\small\thepage}
\renewcommand{\headrulewidth}{0.4pt}

% Elementos da interface: botões, menus, nomes de campos
\newcommand{\ui}[1]{\textsf{\textbf{#1}}}
\newcommand{\menu}[1]{\textsf{#1}}
\newcommand{\sid}[1]{\texttt{#1}}

\newtcolorbox{tip}{colback=tipbg,colframe=accent,boxrule=0.5pt,arc=2pt,
  left=6pt,right=6pt,top=4pt,bottom=4pt,fonttitle=\bfseries,title=Sugestão}
\newtcolorbox{careful}{colback=warnbg,colframe=warnfr,boxrule=0.5pt,arc=2pt,
  left=6pt,right=6pt,top=4pt,bottom=4pt,fonttitle=\bfseries,title=Atenção}

\title{\color{accent}\textbf{Thin Film Logbook}\\[4pt]\large Manual de utilização}
\author{}
\date{\today}

\begin{document}

\maketitle
\thispagestyle{empty}

\begin{abstract}
O Thin Film Logbook é uma aplicação web para registar o que acontece a cada
amostra do laboratório: de onde vem, cada passo de processo (deposição,
recozimento, implantação, \ldots), cada análise (XRD, RBS, SEM, \ldots), e
as fotografias e ficheiros de dados correspondentes. Este manual explica a
utilização do dia a dia. A Secção~\ref{sec:admin} descreve as tarefas
adicionais dos administradores, e a Secção~\ref{sec:install} como manter
o servidor e fazer cópias de segurança.
\end{abstract}

\tableofcontents
\newpage

%=====================================================================
\section{Primeiros passos}
%=====================================================================

\subsection{Abrir o logbook}

O logbook corre num computador do laboratório e abre-se em qualquer
navegador, num computador ou num telemóvel, no endereço indicado pelo
administrador, por exemplo \texttt{http://10.10.13.172:8000}. É preciso
estar ligado à rede do laboratório.

\subsection{A sua conta}

Não há registo livre: é um administrador que cria a sua conta. Recebe um
email com o nome de utilizador e uma palavra-passe temporária (ou o
administrador entrega-lhos).

\begin{enumerate}
  \item Inicie sessão com o nome de utilizador e a palavra-passe temporária.
  \item Escolha uma nova palavra-passe.
  \item Indique o seu \ui{Nome} (aparece como investigador(a) nas suas
        amostras) e as suas \ui{Iniciais}: 2--4 letras, únicas no
        laboratório. As iniciais entram nos IDs das suas amostras, por
        exemplo \sid{ASS-2026-001}.
\end{enumerate}

\begin{careful}
Escolha as iniciais com cuidado. Depois de ter amostras, só um
administrador as pode alterar, e os IDs das amostras já criadas nunca mudam.
\end{careful}

Se se esquecer da palavra-passe, use \ui{Esqueceu-se da palavra-passe?} na
página de início de sessão: recebe por email uma ligação de acesso válida
uma única vez. Pode alterar a palavra-passe e o nome a qualquer momento no
seu perfil (clique no seu nome, no fundo da barra lateral).

\subsection{Orientar-se na aplicação}

A barra lateral à esquerda (um botão de menu no telemóvel) tem:

\begin{center}
\begin{tabularx}{\linewidth}{@{}lX@{}}
\toprule
\menu{Amostras}       & Todas as amostras do laboratório, das mais recentes para as mais antigas. \\
\menu{Pesquisa}       & Filtrar, agrupar, partilhar e exportar amostras. \\
\menu{Substratos}     & A biblioteca partilhada de substratos a partir dos quais se fazem as amostras. \\
\menu{Projetos}       & Projetos de colaboração com outras instituições. \\
\menu{Nova amostra}   & Criar uma amostra. \\
\menu{Sugestões}      & Enviar uma ideia, um problema ou uma pergunta aos administradores. \\
\menu{Caixa de entrada}, \menu{Configurações} & Só para administradores (Secção~\ref{sec:admin}). \\
\bottomrule
\end{tabularx}
\end{center}

No fundo da barra lateral estão o seletor de idioma \ui{PT / EN} e o botão
lua/sol para o modo escuro. Os termos técnicos (nomes de técnicas, unidades,
valores das listas) ficam em inglês nos dois idiomas.

%=====================================================================
\section{Conceitos principais}
%=====================================================================

\begin{center}
\begin{tikzpicture}[
  box/.style={draw=accent,rounded corners=2pt,fill=tipbg,minimum width=2.8cm,
              minimum height=0.8cm,font=\sffamily\small,align=center},
  arr/.style={-{Stealth},thick,accent!70}]
  \node[box] (sub) {Substrato\\[-2pt]\scriptsize biblioteca};
  \node[box,right=1.2cm of sub] (sam) {Amostra\\[-2pt]\scriptsize\texttt{ASS-2026-001}};
  \node[box,left=1.2cm of sub, yshift=0cm] (prj) {Projeto\\[-2pt]\scriptsize colaborações};
  \node[box,below left=1.1cm and -0.6cm of sam] (stp) {Passos de processo};
  \node[box,below right=1.1cm and -0.6cm of sam] (ana) {Análises};
  \node[box,below=0.9cm of stp,xshift=1.9cm,minimum width=3.5cm] (att) {Fotografias e ficheiros};
  \draw[arr] (sub) -- node[above,font=\scriptsize]{feita de} (sam);
  \draw[arr] (prj) to[bend left=25] node[above,font=\scriptsize]{enviada para} (sam);
  \draw[arr] (sam) -- (stp);
  \draw[arr] (sam) -- (ana);
  \draw[arr] (stp) -- (att);
  \draw[arr] (ana) -- (att);
  \draw[arr,dashed] (sam.east) to[out=0,in=0,looseness=2.5]
        node[right,font=\scriptsize,align=left]{dividida em\\pedaços} (sam.north east);
\end{tikzpicture}
\end{center}

\begin{description}
  \item[Amostra] Uma peça física de material, com um ID gerado que nunca
    muda, um responsável (o investigador) e uma \emph{cronologia}.
  \item[Cronologia] A lista ordenada dos passos de processo e análises
    feitos à amostra, cada um com os seus parâmetros, notas, data,
    fotografias e ficheiros.
  \item[Substrato] Uma entrada na biblioteca partilhada (material,
    orientação, espessura, polimento, fornecedor, ficha técnica). As
    amostras de filme fino e bulk são feitas a partir de um substrato.
  \item[Projeto] Uma colaboração com outra universidade ou centro que nos
    envia amostras. O seu código curto substitui as iniciais nos IDs dessas
    amostras.
  \item[Pedaços] Dividir uma amostra cria novas amostras que herdam a sua
    história até ao corte.
\end{description}

\subsection{IDs das amostras}

\begin{center}
\begin{tabular}{@{}ll@{}}
\toprule
\sid{ASS-2026-001}       & Investigador(a) ASS, ano 2026, a sua primeira amostra desse ano \\
\sid{ASS-2026-001-A}     & Pedaço A cortado dessa amostra \\
\sid{ASS-2026-001-A-2}   & Pedaço 2 cortado do pedaço A \\
\sid{ASS-2026-001-REF}   & Pedaço de referência \\
\sid{UMIN1-2026-001}     & Primeira amostra de 2026 do projeto UMIN1 \\
\bottomrule
\end{tabular}
\end{center}

Os IDs são atribuídos quando a amostra é guardada e nunca mudam, mesmo que
a amostra passe para outra pessoa. Use o \ui{Rótulo} (opcional) para a sua
própria descrição (``teste de guia de onda, lote 2''); não precisa de ser
único.

\subsection{Quem vê o quê}

Todas as pessoas com sessão iniciada veem o \emph{resumo} de todas as
amostras: ID, rótulo, tipo, materiais, substrato, etiquetas, técnicas
utilizadas e investigador. Os \emph{detalhes} (descrições, parâmetros,
notas, datas, ficheiros) só são visíveis para:
\begin{itemize}
  \item o responsável pela amostra;
  \item os administradores;
  \item as pessoas com quem o responsável a partilhou (\ui{Ver} ou \ui{Editar});
  \item num passo ou análise em particular, as pessoas indicadas em
        \ui{Realizado por}.
\end{itemize}

As amostras sem responsável estão abertas a todos.

%=====================================================================
\section{Amostras}
%=====================================================================

\subsection{Criar uma amostra}

Clique em \menu{Nova amostra} e escolha o \ui{Tipo de amostra}:

\begin{center}
\begin{tabularx}{\linewidth}{@{}lX@{}}
\toprule
\textbf{Tipo} & \textbf{O que preencher} \\
\midrule
Thin film     & O substrato da biblioteca (obrigatório), as dimensões e, depois
                da deposição, o material do filme. \\
Bulk          & O substrato da biblioteca de que é um pedaço (obrigatório), as
                dimensões, e um \ui{ID externo} se foi cortado por outra pessoa. \\
Commercial    & Descrição, fornecedor, n.º de produto/lote e a ficha técnica
                (PDF ou fotografia, obrigatória). \\
Collaboration & O projeto (obrigatório), uma descrição (obrigatória) e,
                opcionalmente, o ID que o parceiro lhe deu. \\
\bottomrule
\end{tabularx}
\end{center}

Os nomes dos tipos ficam em inglês, como na aplicação. Todos os tipos têm
\ui{Rótulo} e \ui{Etiquetas}. Se o substrato ou o projeto de que precisa
não estiver na lista, siga a ligação \ui{Falta algum?}: abre o formulário
noutro separador; acrescente-o lá e depois recarregue o formulário da
amostra.

O campo \ui{Investigador(a)} é, por omissão, o seu. Escolha outra pessoa
para criar a amostra em nome dela (Secção~\ref{sec:assign}).

\subsection{A página da amostra}

A página de uma amostra mostra o resumo no topo, depois o painel
\ui{Partilha} (para o responsável) e depois a cronologia. A partir daqui
pode:
\begin{itemize}
  \item \ui{Editar} a amostra ou \ui{Eliminar} (só o responsável e os
        administradores; não enquanto tiver pedaços);
  \item \ui{Adicionar passo de processo} e \ui{Adicionar análise};
  \item \ui{Dividir em pedaços} (Secção~\ref{sec:cut});
  \item reordenar a cronologia (Secção~\ref{sec:timeline}).
\end{itemize}

%=====================================================================
\section{Registar passos e análises}
%=====================================================================

\subsection{O formulário}

Escolha primeiro o tipo; o formulário mostra então os campos desse tipo.

\begin{description}
  \item[Data] Opcional, só o dia. O botão \ui{Hoje} preenche-a.
  \item[Onde] CTN, INESC MN ou outro instituto.
  \item[Equipamento] A máquina utilizada. A lista só mostra o equipamento
    que existe no local escolhido em \ui{Onde}. (Numa deposição chama-se
    \ui{Máquina}.)
  \item[Realizado por] Por omissão, o próprio. Acrescente colegas, ou
    \ui{Outra pessoa (escreva o nome)} para alguém sem conta.
  \item[Parâmetros] Os campos próprios do tipo (ver abaixo).
  \item[Outros parâmetros] Tudo o resto, um \texttt{chave: valor} por linha.
  \item[Notas] Texto livre.
  \item[Fotografias e ficheiros] Podem ser adicionados aqui ou mais tarde,
    na cronologia.
\end{description}

\subsection{Preencher os parâmetros}

\begin{itemize}
  \item \textbf{Números} têm a sua caixa e unidade. As pressões estão em
    mTorr e podem escrever-se como \texttt{1e-4}; a vírgula decimal também
    funciona.
  \item \textbf{Durações} aceitam \texttt{min:s} (\texttt{12:30}),
    \texttt{h:min:s} ou só minutos (\texttt{12}).
  \item \textbf{Listas} oferecem os valores já guardados. Escolha
    \ui{+ Adicionar novo\ldots} para escrever um novo, que passa a ser
    oferecido a todos. Veja primeiro a lista, para que o mesmo alvo ou gás
    não fique escrito de três maneiras diferentes.
  \item \textbf{Variantes}: alguns tipos perguntam primeiro qual a variante,
    por exemplo Annealing $\to$ Furnace ou RTA, e mostram depois os campos
    dessa variante.
  \item \textbf{Modos} (XRD / HRXRD): marque cada modo medido
    (2$\theta$/$\omega$, rocking curve, grazing incidence, \ldots) e
    preencha a gama, o passo, o tempo de aquisição e a configuração.
  \item \textbf{Caixas de seleção} mostram campos adicionais, por exemplo
    \ui{Amostra de calibração} numa deposição.
\end{itemize}

\subsection{Tipos disponíveis}

\begin{center}
\begin{tabularx}{\linewidth}{@{}lX@{}}
\toprule
\textbf{Passos de processo} & Deposition, Annealing (Furnace / RTA), Ion
  Implantation, Lithography (optical / e-beam), Etching (dry / wet), Cut,
  Other Process \\
\midrule
\textbf{Análises} & XRD / HRXRD, XPS, SEM, TEM, RBS (RBS / Channeling),
  AFM, Raman, Optical Microscopy, Hall / 4-probe, Ellipsometry,
  Profilometry, Other Analysis \\
\bottomrule
\end{tabularx}
\end{center}

Os administradores podem acrescentar tipos ou alterar os seus campos
(Secção~\ref{sec:admin}).

\subsection{Um ensaio em várias amostras}
\label{sec:runs}

Quando a mesma deposição ou sessão de medidas abrange várias amostras,
registe-a uma só vez: em \ui{Adicionar também a outras amostras}, marque
as outras amostras. Cada uma recebe a sua cópia da entrada, ligadas entre
si como um ensaio (``Também em: \ldots'' na página da amostra).

\begin{itemize}
  \item As fotografias e ficheiros dos campos principais vão para todas as
        amostras (guardados uma só vez).
  \item Ao lado de cada amostra marcada aparece um campo para os ficheiros
        dessa amostra (por exemplo, os seus dados).
  \item Pode acrescentar mais amostras ao ensaio mais tarde, no formulário
        de edição.
  \item Ao editar uma cópia, pode aplicar as alterações às outras. Se isso
        alterar ou apagar um valor que outra amostra já tem, a aplicação
        mostra quais e pede confirmação.
  \item Cada cópia mantém o seu lugar na cronologia da sua amostra, e
        eliminar uma cópia só a remove dessa amostra.
\end{itemize}

\subsection{Fotografias e ficheiros}

Cada cartão da cronologia tem \ui{Adicionar fotografias} e \ui{Anexar
ficheiros}; também pode arrastar ficheiros para um cartão. O envio começa
assim que os ficheiros são escolhidos.

\begin{description}
  \item[Adicionar fotografias] No telemóvel, abre a câmara ou a galeria;
    pode escolher várias de uma vez. As fotografias grandes são endireitadas
    e reduzidas a no máximo 2560\,px / 2\,MB em JPEG; as fotografias HEIC do
    iPhone são convertidas. Os dados de localização são removidos.
  \item[Anexar ficheiros] Guardados exatamente como foram enviados. Use
    esta opção para dados em bruto (por exemplo TIFFs de 16 bits de SEM ou
    AFM, \texttt{.raw}, \texttt{.csv}), para não se perder nada.
\end{description}

As imagens têm uma miniatura na cronologia; clique nela para abrir o
ficheiro completo.

\begin{tip}
Fotografe a amostra, o porta-amostras ou o ecrã do instrumento diretamente
do telemóvel, na bancada: abra a página da amostra, toque em
\ui{Adicionar fotografias} no cartão do passo e tire a fotografia.
\end{tip}

\subsection{A ordem da cronologia}
\label{sec:timeline}

Os passos e as análises partilham a mesma cronologia. Uma nova entrada com
data fica depois de tudo o que tem data igual ou anterior; uma entrada sem
data vai para o fim.

\begin{itemize}
  \item Use as setas para cima/baixo para mover uma entrada, por exemplo
        para ordenar duas coisas feitas no mesmo dia.
  \item \ui{Ordenar por data} repõe as entradas com data por ordem
        cronológica, mantendo a sua ordem dentro de cada dia. As entradas
        sem data ficam onde estão.
\end{itemize}

%=====================================================================
\section{Dividir uma amostra em pedaços}
\label{sec:cut}
%=====================================================================

\ui{Dividir em pedaços}, na página da amostra, regista um passo \emph{Cut}
e cria uma nova amostra por cada pedaço.

\begin{enumerate}
  \item Escolha o \ui{Número de pedaços}. Marque \ui{Guardar também um
        pedaço de referência} para acrescentar mais um (\sid{-REF}).
  \item Tire uma fotografia dos pedaços com a letra de cada um escrita ao lado.
  \item Para cada pedaço, indique opcionalmente um rótulo e as dimensões, e
        escolha o investigador (por omissão, o responsável pela amostra-mãe).
\end{enumerate}

\begin{center}
\begin{tikzpicture}[
  ent/.style={draw=accent,fill=tipbg,rounded corners=1pt,minimum width=2.4cm,
              minimum height=0.5cm,font=\sffamily\scriptsize},
  inh/.style={ent,draw=gray!60,fill=gray!10,text=gray},
  cut/.style={ent,fill=warnbg,draw=warnfr},
  lbl/.style={font=\sffamily\small\bfseries}]
  % amostra-mãe
  \node[lbl] at (0,0.6) {\sid{ASS-2026-001}};
  \node[ent] (p1) at (0,0) {Deposition};
  \node[ent] (p2) at (0,-0.6) {XRD};
  \node[cut] (p3) at (0,-1.2) {Cut};
  \node[ent] (p4) at (0,-1.8) {Annealing};
  % pedaço A
  \node[lbl] at (4.5,0.6) {\sid{ASS-2026-001-A}};
  \node[inh] at (4.5,0) {Deposition};
  \node[inh] at (4.5,-0.6) {XRD};
  \node[inh] (a3) at (4.5,-1.2) {Cut};
  \node[ent] at (4.5,-1.8) {RBS};
  % pedaço B
  \node[lbl] at (8.6,0.6) {\sid{ASS-2026-001-B}};
  \node[inh] at (8.6,0) {Deposition};
  \node[inh] at (8.6,-0.6) {XRD};
  \node[inh] (b3) at (8.6,-1.2) {Cut};
  \node[ent] at (8.6,-1.8) {SEM};
  \draw[-{Stealth},thick,warnfr] (p3.east) -- (a3.west);
  \draw[-{Stealth},thick,warnfr] (p3.east) .. controls (2.3,-3.0) and (6.8,-3.0) .. (b3.west);
  \node[font=\scriptsize,text=gray] at (4.5,-3.0) {cinzento = herdado, só de leitura no pedaço};
\end{tikzpicture}
\end{center}

Cada pedaço:
\begin{itemize}
  \item \textbf{herda} a cronologia da amostra-mãe até ao passo Cut,
        mostrada a cinzento com ``de \ldots''. Nada é copiado: corrigir
        uma entrada na amostra-mãe corrige-a em todos os pedaços. As
        entradas herdadas editam-se na amostra-mãe.
  \item copia o tipo, o substrato, o filme, as etiquetas e a partilha da
        amostra-mãe.
  \item tem as suas próprias entradas depois do corte, separadas das dos
        outros pedaços.
\end{itemize}

A amostra-mãe continua a poder ser usada; o que lhe for acrescentado
depois pertence só a ela. Quem vê um pedaço também vê a história herdada.

%=====================================================================
\section{Partilha e responsabilidade}
%=====================================================================

\subsection{Partilhar uma amostra}

Na página da amostra, no painel \ui{Partilha}, escolha \ui{Partilhar
com\ldots}, a pessoa, e \ui{Ver} ou \ui{Editar}:

\begin{description}
  \item[Ver] Ler os detalhes e descarregar os ficheiros.
  \item[Editar] Também acrescentar e alterar passos, análises e ficheiros.
\end{description}

Para partilhar muitas amostras de uma vez, vá a \menu{Pesquisa}, marque as
amostras (ou \ui{selecionar todas as minhas}) e use \ui{Partilhar as
amostras selecionadas com}. Só o responsável ou um administrador pode
alterar a partilha ou eliminar uma amostra.

\subsection{Realizado por}

Indicar alguém em \ui{Realizado por} permite-lhe ver e editar essa entrada
(com as suas fotografias e ficheiros), mesmo sem acesso ao resto da
amostra. A pessoa é avisada por email. É a forma mais simples de deixar o
técnico que fez o XRD carregar os dados.

\subsection{Passar uma amostra a outra pessoa}
\label{sec:assign}

Pode tornar outro investigador responsável por uma amostra:
\begin{itemize}
  \item no formulário \menu{Nova amostra} (o ID usa então as iniciais dessa
        pessoa);
  \item no formulário \ui{Editar} da amostra;
  \item pedaço a pedaço, no formulário de divisão.
\end{itemize}

Se não for um administrador a fazê-lo, o novo responsável tem de
confirmar. Recebe um email e vê a amostra em \emph{Atribuídas a si ---
confirme, por favor} na página \menu{Amostras}, onde pode \ui{Aceitar} ou
\ui{Recusar} (devolvendo-a). Até lá, e também depois, quem a atribuiu
mantém acesso de edição.

%=====================================================================
\section{Substratos e projetos}
%=====================================================================

\subsection{Biblioteca de substratos}

O separador \menu{Substratos} lista os substratos com material, espessura,
orientação, polimento, tamanho, fornecedor, notas e fichas técnicas.
Qualquer pessoa pode acrescentar um; quem o acrescentou e os
administradores podem editá-lo. Um substrato usado por alguma amostra não
pode ser eliminado.

\subsection{Projetos}

O separador \menu{Projetos} lista as colaborações: nome, instituição
parceira, código, país, contacto e notas. O código (2--6 letras ou
algarismos) é usado nos IDs das amostras do projeto; por isso não pode
coincidir com as iniciais de ninguém, e fica fixo assim que uma amostra o
usa. Uma instituição pode ter vários projetos; escolha-a entre as
sugestões, para ficar sempre escrita da mesma maneira.

%=====================================================================
\section{Pesquisa e exportação}
%=====================================================================

A \menu{Pesquisa} filtra as amostras por texto livre, investigador, tipo de
amostra, material, técnica e etiqueta, e agrupa os resultados (por
investigador, ano, tipo, material, técnica, etiqueta ou projeto). O texto
livre procura em IDs, rótulos, descrições, notas, parâmetros e nomes de
projetos, mas só nas amostras cujos detalhes pode ver.

\ui{Exportar resultados (CSV)} descarrega as amostras encontradas, com uma
linha por passo ou análise e uma coluna \texttt{param:\ldots} por
parâmetro. Abre-se no Excel, no Origin ou em Python
(\texttt{pandas.read\_csv}). Os detalhes ficam em branco nas amostras que
não pode ver.

%=====================================================================
\section{Sugestões}
%=====================================================================

O separador \menu{Sugestões} serve para tudo: um erro, um campo em falta,
um valor de lista que não consegue acrescentar, uma ideia. Fica registada
a página onde estava. Os administradores respondem aí; verá um aviso
quando o fizerem, e pode responder de volta.

%=====================================================================
\section{Para administradores}
\label{sec:admin}
%=====================================================================

As \menu{Configurações} têm três partes.

\begin{description}
  \item[Utilizadores] Criar contas (a palavra-passe temporária é enviada
    por email, ou mostrada no ecrã se o email não estiver configurado),
    repor palavras-passe, tornar ou deixar de ser administrador, alterar o
    nome ou as iniciais de alguém, eliminar utilizadores.
  \item[Tipos de passo de processo / Tipos de análise] Acrescentar, editar
    ou eliminar (quando não estão em uso) tipos. Cada tipo tem um
    \emph{esquema de parâmetros}: uma lista JSON de campos. Editá-lo nunca
    altera os dados já guardados.
  \item[Exportar a base de dados completa (CSV)] Tudo, independentemente
    da partilha.
\end{description}

A \menu{Caixa de entrada} reúne as sugestões de todos; responda, e
marque-as como resolvidas ou reabra-as.

\subsection{Esquemas de parâmetros}

Cada campo tem uma \texttt{key}, um \texttt{label} e um \texttt{type}:

\begin{center}
\small
\begin{tabularx}{\linewidth}{@{}l>{\raggedright}p{4.6cm}X@{}}
\toprule
\textbf{type} & \textbf{preenchido como} & \textbf{chaves adicionais} \\
\midrule
\texttt{number}   & uma caixa numérica & \texttt{unit}, \texttt{placeholder} \\
\texttt{duration} & \texttt{min:s}, \texttt{h:min:s} ou minutos & \texttt{unit}, \texttt{placeholder} \\
\texttt{text}     & uma caixa de uma linha & \texttt{placeholder} \\
\texttt{longtext} & uma caixa grande & \texttt{placeholder} \\
\texttt{choice}   & uma lista partilhada & \texttt{list}; \texttt{located: true} para equipamento \\
\texttt{variant}  & escolher uma opção, cada uma com os seus campos & \texttt{options} \\
\texttt{check}    & uma caixa de seleção que mostra campos & \texttt{fields} \\
\texttt{modes}    & marcar vários; cada um tem os mesmos campos & \texttt{list}, \texttt{fields} \\
\bottomrule
\end{tabularx}
\end{center}

Exemplo (Annealing):

\begin{tcolorbox}[colback=gray!5,colframe=gray!40,boxrule=0.4pt,arc=1pt]
\footnotesize
\begin{verbatim}
[{"key": "method", "label": "Method", "type": "variant", "options": [
   {"value": "Furnace", "fields": [
     {"key": "temperature", "label": "Temperature", "type": "number", "unit": "°C"},
     {"key": "duration", "label": "Duration", "type": "number", "unit": "min"}]},
   {"value": "RTA", "fields": [
     {"key": "atmosphere", "label": "Atmosphere", "type": "choice", "list": "atmosphere"},
     {"key": "temperature", "label": "Temperature", "type": "number", "unit": "°C"},
     {"key": "duration", "label": "Duration", "type": "number", "unit": "s"},
     {"key": "ramp", "label": "Ramp", "type": "number", "unit": "°C/s"}]}]}]
\end{verbatim}
\end{tcolorbox}

Os rótulos (\texttt{label}) escrevem-se em inglês; a tradução portuguesa
fica em \texttt{backend/translations/}.

%=====================================================================
\section{Manter o servidor}
\label{sec:install}
%=====================================================================

Os detalhes completos estão no \texttt{README.md} do repositório. Em resumo:

\begin{tcolorbox}[colback=gray!5,colframe=gray!40,boxrule=0.4pt,arc=1pt]
\footnotesize
\begin{verbatim}
cp .env.example .env          # definir SECRET_KEY, BASE_URL, SMTP_PASSWORD
docker compose up -d
docker compose exec app python -m backend.create_admin
\end{verbatim}
\end{tcolorbox}

\subsection{O que guardar nas cópias de segurança}

\begin{center}
\begin{tabular}{@{}ll@{}}
\toprule
Base de dados (amostras, cronologias, utilizadores) & o volume Docker \texttt{pgdata} \\
Fotografias e ficheiros de dados                    & a pasta \texttt{uploads/} \\
Definições e palavras-passe                         & \texttt{.env} \\
\bottomrule
\end{tabular}
\end{center}

\begin{tcolorbox}[colback=gray!5,colframe=gray!40,boxrule=0.4pt,arc=1pt]
\footnotesize
\begin{verbatim}
docker compose exec -T db pg_dump -U logbook logbook > backup_$(date +%Y%m%d).sql
\end{verbatim}
\end{tcolorbox}

Guarde a cópia da base de dados junto com uma cópia de \texttt{uploads/}.
Essa pasta tem uma subpasta por ID de amostra e, dentro dela, uma por
passo ou análise, com os ficheiros com os nomes originais (por exemplo
\nolinkurl{ASS-2026-001/2026-10-22_XRD-HRXRD/scan01.raw}), mas é a base de
dados que liga cada ficheiro à sua entrada.

\begin{careful}
\texttt{docker compose down} para o logbook e mantém os dados.
\texttt{docker compose down -v} \textbf{apaga a base de dados}.
\end{careful}

\end{document}
